\documentclass[12pt]{article}
\usepackage[utf8]{inputenc}
\usepackage[T1]{fontenc}
\usepackage{hyperref}
\title{Timbe Language Resources for K-12 Curriculum Development}
\author{Compiled by Claude Sonnet 5 (claude-sonnet-5) for Andr\'as Kornai}
\date{}
\begin{document}
\maketitle

\section{Introduction}
Timbe (ISO 639-3 tim; Glottocode timb1251) is a Trans-New Guinea language of Morobe Province, Papua New Guinea. Joshua Project places the Timbe people mainly in the Kabwum District, along the Timbe River valley. Search found only scripture materials. It found no primers, dictionary or grammar.

\section{Primary Educational Resources}
Nothing found in search.

\section{Linguistic and Reference Materials}
Nothing found in search. No Timbe dictionary or grammar appeared in the results.

\section{Religious and Community Resources}
\begin{itemize}
\item Timbe New Testament, \emph{Yusugat Pat Alep Den}, first published in 1987. Its eBible.org page lists a Creative Commons licence (attribution, noncommercial, no derivatives). \url{https://ebible.org/Bible/details.php?id=tim}
\item Companion text on eBible. \url{https://ebible.org/study/content/texts/tim/about.html}
\item PDF of the New Testament on Scripture Earth. \url{https://scriptureearth.org/data/tim/PDF/00-NTtim-web.pdf}
\item Printable PDF (A4) on eBible. \url{https://eBible.org/pdf/tim/tim_a4.pdf}
\item Joshua Project profile of the Timbe people. \url{https://www.joshuaproject.net/people_groups/15491}
\end{itemize}

\section{Digital and Audio Resources}
No audio found in search.

\section{Curriculum Development Context}
The New Testament is the only substantial written text found. It could serve as a source for orthography, reading practice and a text corpus. Because no dictionary or grammar turned up, curriculum work would need new documentation with the community.
\end{document}